\section{Provenance: what came from where}
\label{app:provenance}

This document is a merge of nine independently prepared proposals for the same
system. All nine are retained verbatim under \code{proposals/} so that every
claim here can be audited against its source. Almost nothing in this article is
new mathematics: each statement is carried over from at least one of the nine,
or from the two external sources named in \S\ref{sec:external-sources}, and
this appendix records which.

\subsection{The nine sources}

\begin{longtable}{@{}>{\raggedright\arraybackslash}p{0.10\textwidth}>{\raggedright\arraybackslash}p{0.38\textwidth}>{\raggedright\arraybackslash}p{0.46\textwidth}@{}}
\toprule
Label & Directory, and the archive name it arrived under & Distinguishing emphasis \\
\midrule
\endhead
\src{p1} & \code{p1-structural-search}\newline
\emph{from} \code{FORGE\_Beyond\_Grind(1) (1)/} &
Six certificate families under one measurement protocol; the only
zero-preserving univariate worker; the hardened certificate decoder. \\
\src{p2} & \code{p2-obligation-controller}\newline
\emph{from} \code{Forge-design-and-prototypes/} &
The three-store architecture and the demand protocol; Farkas affine witnesses as
a single feasibility problem; theory-constrained instantiation. \\
\src{p3} & \code{p3-planner-certificate-layer}\newline
\emph{from} \code{Forge\_Beyond\_Grind/} &
Two graphs rather than one fact bag; compositional acceptance; the quantitative
Bernstein obstruction; the Lean runtime design contract. \\
\src{p4} & \code{p4-theory-cooperation}\newline
\emph{from} \code{Forge\_Beyond\_Grind(1)/} &
The largest artefact corpus; residue witnesses; the named \grind{}
solver-extension API; the $2\times2$ structural ablation. \\
\src{p5} & \code{p5-certificate-first}\newline
\emph{from} \code{Forge\_Tactic\_Design/} &
Target-driven square proposal; tabling against cyclic pseudo-proofs; solver
islands and translation ledgers; expected-outcome benchmarking. \\
\src{p6} & \code{p6-proof-logging-cdcl}\newline
\emph{from} \code{forge-beyond-grind/} &
Proof-logging CDCL(T) with independent RUP replay; the integer-negation trap;
the \code{nativeEqTrue} trust finding. \\
\src{p7} & \code{p7-successor-architecture}\newline
\emph{from} \code{forge-grind/} &
Integer-lattice witnesses and certificates of impossibility; the evidence
broker; conclusion indexing; metavariable ownership. \\
\src{p8} & \code{p8-obligation-broker}\newline
\emph{from} \code{forge-tactic-design-and-prototypes/} &
The cleanest dictionary ablation; certificate cost as a scheduling term; the
largest Horn workload; checkers that take the target externally. \\
\src{p9} & \code{p9-proof-planner}\newline
\emph{from} \code{forge\_beyond\_grind (1)/} &
The executive assessment; variables absent from the conclusion; witness as a
program with an obligation; the rank-one dictionary-gap argument. \\
\bottomrule
\end{longtable}

\subsection{Where each section's material came from}

\begin{longtable}{@{}>{\raggedright\arraybackslash}p{0.30\textwidth}>{\raggedright\arraybackslash}p{0.64\textwidth}@{}}
\toprule
Section & Principal sources \\
\midrule
\endhead
\S1 Objective &
The two modes and the preservation proposition appear in all nine; the
four evidence levels are \src{p1}'s framing with \src{p4}'s vocabulary and
\src{p8}'s ``implemented / mathematically justified / proposed / kernel
checked'' distinction. \\
\S2 Baseline &
Five proposals built partial capability tables; the bit-vector and
conditional-rewriting rows are \src{p2}'s, the configuration defaults
\src{p9}'s, the dependency map \src{p2}'s extended with \src{p6}'s and
\src{p9}'s library findings. \\
\S3 Architecture &
Three stores and the demand vocabulary from \src{p2}; two graphs, hyperedges and
scope-as-identity from \src{p3}; the evidence broker, event subscriptions and
metavariable ownership from \src{p7}; the ledger invariant and
certificate-cost-in-benefit from \src{p8}; the nine-tuple planner state and the
six-constructor message vocabulary from \src{p9}; the scheduler score and
reserved quotas from \src{p1}. \\
\S4 Trust &
Three layers and the proof-cost budget from \src{p1}; the five invariants and
the review checklist from \src{p2}; ``a type is not a check'' from \src{p3};
axiom auditing against the pinned version from \src{p4}; the two trust profiles
from \src{p5}; \code{nativeEqTrue} and the insufficiency of a single-name
blacklist from \src{p6}; the four-step acceptance procedure from \src{p9}. \\
\S5 Structural search &
Conclusion indexing and \code{matchDemand} from \src{p7}; constrained matching
from \src{p2}; tabling and the relative-completeness proposition from \src{p5};
the slicing propositions from \src{p3} and \src{p7}; the accumulator obstruction
from \src{p2} and \src{p3}; the flexible-variable IH discipline from \src{p4};
the CEGIS loop and the overfitting trap from \src{p2}; fragment completeness
from \src{p2} and \src{p7}; the tag-variable trick from \src{p7}; the
state-machine invariant format from \src{p5}; Farkas witnesses from \src{p2};
lattice elimination from \src{p7}; residue witnesses from \src{p4}; polynomial
witnesses from \src{p9}; \code{energyAcc} from \src{p3}. \\
\S6 Nonlinear &
The cone language and its soundness from all nine; the quadratic theorem from
\src{p1} and \src{p3}; repair-not-rounding from \src{p1}; target-driven squares
from \src{p5}; the four dictionary-gap arguments from \src{p2}, \src{p5},
\src{p6}, \src{p7} and \src{p9}; the Bernstein lemma and the interior-zero
theorem from \src{p1}, with \src{p3}'s quantitative form; the completeness
sketch from \src{p2} and \src{p3}; square-factor and Sturm certificates from
\src{p1} alone; product envelopes from \src{p1}; the reuse-the-existing-backend
argument from \src{p9} and \src{p5}. \\
\S7 Boolean and external &
The SAT/theory loop from \src{p7}; CDCL(T), RUP replay and the negation trap
from \src{p6}; the resolution-DAG format, \code{verify\_sat} and the exhaustive
oracle from \src{p4}; extraction cost vectors from \src{p7}; solver islands and
translation ledgers from \src{p5}; the reconstruction preference order from
\src{p7}; the no-blanket-combination warning from \src{p5} and \src{p9}. \\
\S8 Integration &
\code{Params}/\code{assertExtra}/\code{mkGoalCore} from \src{p2} and \src{p9};
the solver-extension callbacks from \src{p4}; the rollback-is-not-enough point
from \src{p6}; three reconstruction levels from \src{p2}; the module layout
merged from \src{p1}, \src{p2}, \src{p3}, \src{p5} and \src{p9}; the acceptance
contract table from \src{p5}; the deployment pattern from \src{p9}; the gates
merged from all nine. \\
\S9 Experiments & Each run reports itself. \\
\S10 Evaluation &
Leakage control from \src{p1}, \src{p5} and \src{p9}; the three input conditions
from \src{p2}; the two tracks from \src{p4}; the non-communicating portfolio
from \src{p1} and \src{p6}; metric discipline from \src{p4} and \src{p9}; the
illustrative acceptance target from \src{p9}. \\
\S A Schemas &
\src{p5}'s field names as the base, \src{p3}'s \code{kind} discriminator,
\src{p8}'s \code{\{index, multiplier\}} ideals, \src{p1}'s bounded decoder,
\src{p3}'s duplicate-key rejection, \src{p6}'s anti-vacuity guards and
ignore-the-residual rule, \src{p9}'s typed decoding gate. \\
\bottomrule
\end{longtable}

\subsection{Ideas that survive in exactly one source}

These were the merge's first concern, since a careless consolidation loses
precisely the material that distinguishes one submission from the rest.

\begin{itemize}[leftmargin=1.4em]
\item \textbf{Square-factor and Sturm certificates} (\src{p1}) --- the only
  mechanism anywhere in the collection that can certify a polynomial with an
  interior zero, and the constructive answer to Theorem~\ref{thm:zero}.
\item \textbf{Integer-lattice elimination with a divisibility obstruction}
  (\src{p7}) --- the only worker that emits a certificate of \emph{infeasibility}.
\item \textbf{Finite-residue witnesses} (\src{p4}) --- the only
  existential-over-$\Z$-by-case-analysis certificate.
\item \textbf{Polynomial witnesses discharged by two cone certificates}
  (\src{p9}) --- the smallest complete instance of the composition pattern the
  whole architecture argues for.
\item \textbf{Farkas affine witnesses as one feasibility problem} (\src{p2}) ---
  the observation that synthesis and proof search are jointly linear.
\item \textbf{Proof-logging CDCL(T) with independent RUP re-derivation}
  (\src{p6}), and the resolution-DAG alternative (\src{p4}).
\item \textbf{Theory-constrained quantifier instantiation} (\src{p2}) ---
  emitting an equality constraint at an interpreted subterm instead of failing.
\item \textbf{Certificate cost inside the scheduling benefit} (\src{p8}).
\item \textbf{The quantitative interior-zero coefficient}
  $(\ell-1/3)(u-1/3)<0$ (\src{p3}).
\item \textbf{The rank-one Gram argument} for $(x+y+1)^2$ (\src{p9}) and the
  \textbf{weight-6-versus-coefficient-4 argument} for $(2x-3y)^2$ (\src{p7}).
\item \textbf{The Lean design contracts} --- the runtime and scheduler contract
  (\src{p3}) and the \code{CertificateSpec} soundness contract (\src{p6}); these
  complement rather than duplicate each other, and both now elaborate.
\end{itemize}

\subsection{Disagreements, and how they were resolved}

The nine did not agree on everything. Where they conflicted, this document took
the more conservative position and says so.

\begin{longtable}{@{}>{\raggedright\arraybackslash}p{0.27\textwidth}>{\raggedright\arraybackslash}p{0.67\textwidth}@{}}
\toprule
Disagreement & Resolution \\
\midrule
\endhead
Build or reuse the nonlinear backend &
\src{p9} and \src{p5} identified an existing Lean \code{sos} project and a Lean
\code{lp} project; the other seven proposed building cone and witness machinery
from scratch without mentioning either. Resolved in favour of the conservative
reading: the in-house searches are a cheap first stage, the existing libraries
are the production backend, and ``add an SOS tactic to Lean'' is not claimed as
a contribution (Remark~\ref{rem:sos-reuse}). \\
What \grind{} already does &
\src{p6} conceded conclusion-selected backward patterns to the baseline;
\src{p1} listed them as new work. The concession was adopted. \\
Native-evaluation trust &
Three confidence levels were recorded: an RFC-only reading, a
tactic-reference reading, and \src{p6}'s inspection of \code{Lean.Meta.Native}
and \code{nativeEqTrue}. The most specific source was used, with its own
instruction that the behaviour must be re-audited against the pinned
implementation rather than inherited from a design thread. \\
Is a stable snapshot API available? &
\src{p4} called the solver-extension callbacks ``real integration points, not
hypothetical hooks''; \src{p3} declined to assume any stable public API. Both
are reported: the inventory is \src{p4}'s, the framing is \src{p3}'s. \\
Bernstein split policy &
Round-robin axis bisection (\src{p4}) against largest-relative-width (\src{p3})
against \code{depth mod n} (\src{p2}). All three make every coordinate shrink
along a branch, which is what Proposition~\ref{prop:bernstein-complete} needs;
largest-\emph{absolute}-width does not, and is rejected. \\
Horn forward-search growth &
\src{p5} reports $2^{d+1}-1$ facts and \src{p9} reports $2^d$ on the same rule
family. The difference is entirely the forward baseline's stopping rule. Both
tables are printed, side by side, as a methodological example
(Table~\ref{tab:horn-four}). \\
Quadratic matrix convention &
\src{p3} places a linear coefficient at $H_{0i}=b_i/2$; \src{p4} states it as
$2Q_{0i}$. These are inverse readings of the same matrix. \S\ref{sec:quadratic}
uses one convention throughout. \\
Accumulator naming &
\code{ra}/\code{la} (\src{p3}), \code{qrev} (\src{p4}), \code{revAux}
(\src{p2}), \code{revAcc} (\src{p9}). One naming is used per derivation; they
are never mixed inside one chain of equalities. \\
Pigeonhole encodings &
\src{p4} and \src{p6} use different clause sets and report different columns.
Both tables are kept, explicitly labelled as non-comparable. \\
\bottomrule
\end{longtable}

\subsection{Material from outside the nine}
\label{sec:external-sources}

Two neighbouring projects were reviewed after the merge, and a substantial
amount of this document comes from them rather than from any proposal. They
matter because between them they had \emph{implemented} the thing all nine
proposals specified and none built: a boundary at which nothing becomes a fact
without a checked Lean proof.

\begin{itemize}[leftmargin=1.4em]
\item \textbf{Djex} --- a type-directed Haskell expression synthesizer merging
  Djinn (Dyckhoff's contraction-free \textsf{LJT}, terminating, able to decide
  uninhabitation) with Exference (ranked best-first search with typeclass
  evidence resolution).
\item \textbf{Leant} --- a Djex-powered synthesis REPL for Lean~4, in which
  every candidate term is re-elaborated by Lean before the user sees it.
\end{itemize}

What came from them:

\begin{longtable}{@{}>{\raggedright\arraybackslash}p{0.34\textwidth}>{\raggedright\arraybackslash}p{0.60\textwidth}@{}}
\toprule
Section & Material \\
\midrule
\endhead
\S\ref{sec:tactic-gap} &
The neighbouring-tactics table. The nine named Aesop, Duper, \code{linarith}
and the SOS and LP projects, but omitted \code{exact?}/\code{apply?},
\code{tauto}/\code{itauto}, \code{decide} and \code{omega}. \\
\S\ref{sec:architecture} &
Deduplication never refunds budget; a branch dropped for a resource reason
stays structurally visible to the top. \\
\S\ref{sec:kernel-traps} &
\code{debug.skipKernelTC} and the asynchronous \code{addDecl} path
(Invariant~\ref{inv:sync}); the unforgeable stamp; the verification pipeline. \\
\S\ref{sec:lattice} &
The status lattice and its three governing invariants, including the
approximation-polarity rule for pruning. \\
\S\ref{sec:boundary-impl} &
The whole subsection: ephemeral \code{example}, \code{noncomputable},
\code{autoImplicit} off, re-elaborating the original goal text, the committed
environment as positive evidence (Invariant~\ref{inv:positive}), the closed
rejection vocabulary, transport-versus-semantic failure, the two-tier gate,
the abstract receipt, and ranking as a sealed capability. \\
\S\ref{sec:negative} &
The whole subsection: exhausted versus truncated; the completeness ledger as a
per-node safety bit computed by the encoder; the three-conjunct gate and its
self-reference sharpening with the monotonicity rule; portfolio asymmetry;
\textsf{LJT}; Glivenko's classical fallback and the rule that quantified goals
skip it. \\
\S\ref{sec:boolean} &
Instance obligations during synthesis: one function for pruning and checking,
termination work at admission time, budget-exhausted as a distinct
constructor, receipts bound to their environment, and class obligations as
structurally disjoint propositions. \\
\S\ref{sec:worker-split} &
The tactic-as-interface, worker-as-executor split and its four costs. \\
\S\ref{sec:identity} &
Three levels of identity; the origin table never leaving the worker
(Invariant~\ref{inv:origin}); environment generation tokens. \\
\S\ref{sec:search-ir} &
Why the search IR should not be raw \code{Expr}; deduplication discarding the
authority-bearing representative. \\
\bottomrule
\end{longtable}

Two qualifications on that material.

The \code{debug.skipKernelTC} and asynchronous-\code{addDecl} findings were
\emph{verified first-hand} against the installed \code{v4.34.0} source rather
than taken on the reviewing project's word, because its own analysis is written
against \code{v4.33.0}. Both hold at the pinned version.

And one claim was corrected in the other direction. An \textsf{LJT} decision
procedure decides non-inhabitation, but the reference implementation emits for
a negative answer a single enumeration tag and no artefact --- no closed-branch
tree, nothing a second program can re-check. It is therefore \emph{weaker} than
lattice elimination, not equivalent to it, and \S\ref{sec:negative} says so.
The opportunity that follows --- reifying a finite failed search tree into a
checkable refutation object --- appears to be unbuilt anywhere, and would be a
genuine contribution rather than a port.

\subsection{What the merge changed}

Four things in this document are not simply inherited.

First, the nine experiment sections are reported as nine, with the non-pooling
rule stated explicitly and the four Horn setups tabulated against each other.
Several of the source articles warned against pooling their \emph{own} numbers;
none could warn against pooling across articles, because none knew about the
others.

Second, the bibliography is deduplicated: roughly 215 \code{bibitem} entries
under about 100 keys became 56 entries under descriptive keys, keeping the
richest description of each work and preserving details only one source had
recorded --- inspected blob SHA-1s, the release timestamp, the
blob-id-is-not-a-commit-id caveat, and the \code{sos} toolchain version
mismatch.

Third, the build is plain pdf\LaTeX{} with stock packages. The nine used four
engines and nine font families between them, and every one of their READMEs had
to explain what to do when a font was missing.

Fourth, and the only new evidence: three Lean files were compiled
(Table~\ref{tab:lean}). Every source recorded \code{NOT\_RUN}. That is now
narrower --- three of sixteen files elaborate under Lean v4.34.0 --- and the
limits of what it shows are stated where the result appears.

\subsection{Conventions used in this document}

\code{\textbackslash status\{...\}} banners, inherited from several of the
sources, mark whether a mechanism was implemented and executed, specified only,
or proposed for production. \src{pN} labels attribute a result to the run that
produced it. \unknown{} always means a bounded search found nothing; it never
means a statement is false.
